\begin{frame}{DSU}
    \metroset{block=fill}

    Podemos implementar uma estrutura de dados melhor, sem a necessidade da utilização de sets. É uma estrutura de dados que chamamos de DSU (Disjoint set union).

    {\bf Ideia principal.} No lugar de manter vários sets, cada conjunto terá um único representante. Isto é implementado de forma que cada nó tenha um pai, e o representante seja a raiz da árvore gerada. Assim, as operações são:
    
    \begin{itemize}
        \item União: sejam $a$ o representante do menor conjunto e $b$ o representante do maior conjunto. Definiremos o novo pai de $a$ sendo $b$.
        \item Procura: para saber se $a$ e $b$ estão no mesmo conjunto, basta ver se possuem o mesmo representante.
    \end{itemize}

    Pela heurística de \textit{small to large}, a partir de cada elemento, é necessário visitar até $\log n$ vértices até chegar na raiz da árvore.


    % \begin{block}{Implementação da busca binária}
    %     \footnotesize
    %     \inputminted{cpp}{codigos/busca.binaria.cpp}    
    % \end{block}

\end{frame}
